% --- Archon physics formalization source begin ---
% archon:physics
% archon:covers IPhO2026Problems/problem_ipho_2026_t3_c3.lean
% archon:source-report reports/ipho_2026/problem_ipho_2026_t3_c3.source.json
% archon:problem-id ipho_2026_t3
% archon:part-id T3-C3
% archon:previous-part ipho_2026_t3_b1
% archon:previous-part-policy derive_inline_from_problem_only_material
% archon:previous-part ipho_2026_t3_c2
% archon:previous-part-policy derive_inline_from_problem_only_material

\chapter{Ipho Physics problem ipho\_2026\_t3\_c3}
\label{ch:IPhO2026Problems_problem_ipho_2026_t3_c3}

\paragraph{Problem source.}
\#\# Physical scenario

The paramagnetic torus executes the Carnot refrigeration cycle
1 -> 2 -> 3 -> 4 -> 1 shown in Figure 3b in the H-versus-T plane.  T\_h and T\_c
are the hot- and cold-reservoir temperatures; Q\_h is the magnitude of heat
delivered to the hot reservoir and Q\_c is the magnitude absorbed from the cold
reservoir.  The equation of state is T*M*V = n*K*H and the isothermal heat
relation from part B may be reused.

\#\# Current subquestion T3-C3

Using the supplied potassium-chromate and liquid-helium data, find the helium temperature after one cycle.

\paragraph{Shared problem context.}
The paramagnetic torus executes the Carnot refrigeration cycle
1 -> 2 -> 3 -> 4 -> 1 shown in Figure 3b in the H-versus-T plane.  T\_h and T\_c
are the hot- and cold-reservoir temperatures; Q\_h is the magnitude of heat
delivered to the hot reservoir and Q\_c is the magnitude absorbed from the cold
reservoir.  The equation of state is T*M*V = n*K*H and the isothermal heat
relation from part B may be reused.

\paragraph{Current subquestion.}
Using the supplied potassium-chromate and liquid-helium data, find the helium temperature after one cycle.

\paragraph{Figure/image path.}
ipho\_2026\_source/image/T3\_page-4.png

\paragraph{Reusable previous-part conclusions.}
\begin{itemize}
\item Source T3-B1. Question: At fixed temperature T, H changes from H\_i to H\_f. Find the heat Q transferred into the torus. Policy: derive\_inline\_from\_problem\_only\_material
\item Source T3-C2. Question: Express M\_1 in terms of M\_2, M\_3, and M\_4. Policy: derive\_inline\_from\_problem\_only\_material
\end{itemize}

\paragraph{Formalization target.}
create a compiling Lean file with sorry bodies at `IPhO2026Problems/problem\_ipho\_2026\_t3\_c3.lean`.
The Lean declarations must preserve the physical quantities, dimensions or dimensional roles, figure labels, governing-law hypotheses, and final relation expressed by this problem.
Use Mathlib/Physlib names found through LeanExplore where available. If a domain API is missing, introduce faithful local abstractions rather than scalar placeholder aliases.
This is an answer-blind solve-phase task. Derive a candidate result only from the problem-side evidence above; do not assume a target value, select a tolerance around a desired value, or encode a desired result into a candidate domain.
For numerical questions, define the raw end-to-end quantity and apply an explicit source-derived rounding rule. For identification questions, characterize or prove uniqueness before recording the derived candidate.

\begin{theorem}[Ipho Physics formalization target]
\label{thm:physics:ipho_2026_t3_c3:target}
\leanok
The assigned autoformalize agent should translate this IPhO physics problem into Lean declarations in the covered file, with theorem and lemma proof bodies written as `by sorry`.
\end{theorem}
\begin{proof}
\leanok
Fully proved in \texttt{IPhO2026Problems/problem\_ipho\_2026\_t3\_c3.lean}
(namespace \texttt{IPhO2026.T3C3}): the cold-leg heat
$Q_c = (\mu_0 n K/(2 T_0))(H_2^2 - H_3^2)$ (\texttt{heat\_absorbed\_on\_cold\_leg},
positive by \texttt{cold\_leg\_heat\_pos}) is converted by constant-$c$
calorimetry (\texttt{helium\_temperature\_drop}) into
$T_f = T_0 - (\mu_0 n K/(2 T_0))(H_2^2 - H_3^2)/(\rho_{He} V_{He} c)$
(\texttt{helium\_final\_temperature}); at the supplied data this evaluates to
$\approx 0.99005\,\mathrm{K} \in (0.989, 0.991)$
(\texttt{rawFinalHeliumTemperature\_bounds} via \texttt{Real.pi\_gt\_d4},
\texttt{Real.pi\_lt\_d4}), which the source-derived rounding rule (nearest
$0.01\,\mathrm{K}$) reports as $0.99\,\mathrm{K}$
(\texttt{reportedHeliumTemperature\_value}).
\end{proof}
% --- Archon physics formalization source end ---
